\subsection{Vertex-boundary formulations}
\label{sec:area_law_vertex_boundaries}

\begin{definition}[Inner boundary and boundary endpoints]
    \label{def:area_law_vertex_boundaries}
    \lean{TNLean.PEPS.AreaLaw.innerBoundary,
        TNLean.PEPS.AreaLaw.endpointBoundary}
    \leanok
    \uses{def:area_law_domain,def:area_law_hamiltonian}
    For a finite induced domain $\Lambda\subset\mathbb Z^2$ and $A\subseteq\Lambda$,
    define
    \begin{align}
        \partial_{\rm in}A
            & =\{x\in A:\exists y\in\Lambda\setminus A,
        \ \{x,y\}\in\partial_\Lambda A\},
        \label{eq:area_vertex_inner}                    \\
        Z_A & =\bigcup_{e\in\partial_\Lambda A}e.
        \label{eq:area_vertex_endpoints}
    \end{align}
    These are the conventions of \cite[Corollary~1.2]{OpenAI2026AreaLaw}.
\end{definition}

\begin{theorem}[Boundaries of the empty and full cuts]
    \label{thm:area_law_vertex_boundaries_trivial}
    \lean{TNLean.PEPS.AreaLaw.innerBoundary_empty,
        TNLean.PEPS.AreaLaw.innerBoundary_univ,
        TNLean.PEPS.AreaLaw.endpointBoundary_empty,
        TNLean.PEPS.AreaLaw.endpointBoundary_univ}
    \leanok
    \uses{def:area_law_vertex_boundaries}
    Both $\partial_{\rm in}A$ and $Z_A$ are empty when $A=\varnothing$ or
    $A=\Lambda$.
\end{theorem}
\begin{proof}\leanok
    In both cases $\partial_\Lambda A=\varnothing$, since no edge joins a site of
    $A$ to a site of $\Lambda\setminus A$.
\end{proof}

\begin{theorem}[Degree bound for induced lattice domains]
    \label{thm:area_law_lattice_degree}
    \lean{TNLean.PEPS.AreaLaw.domainGraph_degree_le_four}
    \leanok
    \uses{def:area_law_domain}
    Every site of a finite induced square-lattice domain has degree at most four.
\end{theorem}
\begin{proof}\leanok
    The neighbors of $(a,b)$ form a subset of
    $\{(a+1,b),(a-1,b),(a,b+1),(a,b-1)\}$.
\end{proof}

\begin{theorem}[Comparison of edge and vertex boundaries]
    \label{thm:area_law_vertex_boundary_counts}
    \lean{TNLean.PEPS.AreaLaw.edgeBoundary_card_le_four_mul_innerBoundary_card,
        TNLean.PEPS.AreaLaw.edgeBoundary_card_le_two_mul_endpointBoundary_card}
    \leanok
    \uses{def:area_law_vertex_boundaries}
    For every cut in a finite induced lattice domain,
    $|\partial_\Lambda A|\le4|\partial_{\rm in}A|$ and
    $|\partial_\Lambda A|\le2|Z_A|$.
    No connectedness assumption is required.
\end{theorem}
\begin{proof}\leanok
    \uses{def:area_law_vertex_boundaries,thm:area_law_lattice_degree}
    Let $G_A$ have exactly the crossing edges, with degrees $d_A(x)$.
    It is bipartite between the inner boundaries of $A$ and its complement,
    and its nonisolated sites are exactly $Z_A$.
    Counting degrees on one side and on both sides gives
    \begin{align}
        |\partial_\Lambda A|
         & =\sum_{x\in\partial_{\rm in}A}d_A(x)
        \le4|\partial_{\rm in}A|,               \\
        2|\partial_\Lambda A|
         & =\sum_{x\in Z_A}d_A(x)\le4|Z_A|.
    \end{align}
    Empty boundaries satisfy the same identities.
\end{proof}

\begin{theorem}[Vertex bounds assuming the uniform edge area law]
    \label{thm:area_law_conditional_vertex_bounds}
    \lean{TNLean.PEPS.AreaLaw.regionalEntropy_le_four_mul_innerBoundary_card,
        TNLean.PEPS.AreaLaw.regionalEntropy_le_two_mul_endpointBoundary_card,
        TNLean.PEPS.AreaLaw.UniformAreaLaw.vertex_boundary_bounds}
    \leanok
    \uses{def:area_law_vertex_boundaries,def:area_law_ground_entropy}
    Assume that, for fixed $q,R,J,\Delta$, a nonnegative constant $C$ gives
    $S_\Omega(A)\le C|\partial_\Lambda A|$ uniformly for the finite-domain
    Hamiltonians and ground vectors of Section~\ref{sec:finite_domain_area_law}.
    Then the same constant gives
    $S_\Omega(A)\le4C|\partial_{\rm in}A|$ and
    $S_\Omega(A)\le2C|Z_A|$ uniformly for those instances.
    The two implications also hold for a single entropy inequality.
\end{theorem}
\begin{proof}\leanok
    \uses{thm:area_law_vertex_boundary_counts}
    Multiply the two boundary comparisons by $C\ge0$ and apply the assumed
    edge-boundary entropy inequality. The constant is chosen before the
    domain, Hamiltonian, ground vector, and cut.
\end{proof}
